\section{Aplicacion empresarial}

El analisis financiero permite convertir resultados numericos en decisiones concretas. Si las ventas crecen, pero el costo de ventas aumenta a mayor velocidad, la empresa debe revisar precios, proveedores, procesos productivos y eficiencia operativa. Si el patrimonio pierde participacion relativa frente al total activo, tambien conviene evaluar la estructura de financiamiento.

En la gestion empresarial, el analisis horizontal ayuda a detectar tendencias, mientras que el analisis vertical permite observar la composicion economica y financiera. Ambos metodos pueden usarse en presupuestos, control de costos, evaluacion de rentabilidad, decisiones de inversion y seguimiento de metas.

\section{Aplicacion del analisis financiero en Ingenieria de Sistemas}

La Contabilidad Empresarial y la Ingenieria de Sistemas se relacionan mediante los sistemas de informacion que registran, procesan y reportan datos economicos. En una empresa, las operaciones de ventas, compras, bancos, inventarios y contabilidad pueden integrarse en un ERP, permitiendo que la informacion contable sea consultada y transformada en reportes financieros.

Las bases de datos permiten almacenar ventas, costos, activos, pasivos, patrimonio, compras e inventarios de manera estructurada. A partir de estas tablas, un analista o ingeniero de sistemas puede construir consultas que agrupen informacion por periodo, cuenta contable, area, proyecto o centro de costo.

El analisis horizontal puede automatizarse mediante hojas de calculo, SQL, Python, Power BI o modulos de ERP que comparen el periodo actual con el periodo anterior. Del mismo modo, el analisis vertical puede calcularse dividiendo automaticamente cada partida entre una base comun, como ventas netas o total de activos.

Los dashboards financieros permiten presentar resultados de forma visual, mostrando variaciones, margenes, estructura de costos, liquidez, deuda y rentabilidad. Estos tableros pueden incluir alertas cuando los costos crecen mas rapido que las ventas, cuando el margen bruto disminuye o cuando el endeudamiento supera un limite definido.

Como ejemplo laboral, una empresa tecnologica puede monitorear ventas por licencias, costos de infraestructura cloud, gastos administrativos, soporte tecnico y utilidad neta. Si el gasto cloud pasa de S/ 8,000.00 a S/ 12,000.00, el analisis horizontal mide el incremento; si ese gasto representa un porcentaje creciente de las ventas, el analisis vertical ayuda a evaluar si la estructura de costos sigue siendo sostenible.

Por ello, el analisis financiero no solo es una herramienta contable. Tambien es un campo de aplicacion para la Ingenieria de Sistemas, porque permite disenar reportes, automatizar indicadores y apoyar decisiones gerenciales con informacion oportuna, trazable y util.

\subsection{Aplicacion tecnologica del analisis financiero}

Una solucion tecnologica de analisis financiero puede organizarse como un flujo: ERP, base de datos, proceso ETL, calculo de indicadores, dashboard y alerta gerencial. En este flujo, el ERP registra operaciones; la base de datos almacena transacciones; el ETL limpia y organiza la informacion; el calculo financiero genera variaciones, porcentajes y ratios; y el dashboard presenta resultados visuales para la toma de decisiones.

Herramientas como Excel, SQL, Python, Power BI o sistemas web permiten automatizar reportes y reducir errores de calculo. Esta automatizacion facilita que la gerencia revise tendencias, margenes, deuda, liquidez y rentabilidad con informacion actualizada.

\section{Conclusiones}

\begin{itemize}
  \item El analisis financiero convierte estados financieros en informacion para decidir.
  \item El analisis horizontal permite observar tendencias y variaciones entre periodos.
  \item El analisis vertical permite estudiar la estructura porcentual de un estado financiero.
  \item En la Empresa ABC, las ventas crecen, pero el costo de ventas crece mas rapido, reduciendo el margen bruto.
  \item Los resultados deben complementarse con ratios de liquidez, gestion, solvencia y rentabilidad.
  \item Para Ingenieria de Sistemas, estos analisis pueden automatizarse mediante ERP, bases de datos y dashboards.
\end{itemize}
